% ---------------------------------------------------------------------
\section{Mettre en place l’IAD, étape par étape}
% ---------------------------------------------------------------------

\subsection{Définir sa stratégie de diffusion}

Avant de choisir une plateforme, répondez à une question :
\textbf{où vos lecteurs achètent-ils leurs livres ?}

\begin{center}
\small
\begin{tabularx}{\linewidth}{P{4.6cm}L}
\entete \textbf{Votre objectif} & \textbf{Montage conseillé} \\
Vendre surtout en ligne, budget minimal & Amazon KDP seul, avec vos
  propres ISBN pour garder la main \\
Être commandable en librairie et en ligne & une plateforme distribuée
  via Dilicom (BoD, Bookelis…) ; éventuellement KDP en plus pour Amazon \\
Petite maison avec un catalogue & vos ISBN, une plateforme d’IAD
  « éditeur » (BoD, IngramSpark) pour le fonds, des tirages courts pour
  les nouveautés et les salons \\
Vendre en direct (salons, site, conférences) & petit tirage numérique +
  IAD pour les ventes à distance \\
\bottomrule
\end{tabularx}
\end{center}

\begin{attention}[Un ISBN, un circuit de distribution]
Vous pouvez utiliser le même ISBN sur Amazon KDP et sur une autre
plateforme, à condition que l’ISBN vous appartienne. En revanche, pour
les libraires, \textbf{un même ISBN ne doit avoir qu’un seul
distributeur référencé dans Dilicom}, sinon les commandes se perdent.
Choisissez clairement « qui distribue quoi ».
\end{attention}

\avoir{BTQuLZikE0g}{\textbf{Créer une maison d’édition sans subvention en
France} : le témoignage d’Ephyra Éditions sur le lancement d’une maison
indépendante.}

\subsection{Préparer le livre}

\begin{checklist}
  \item Manuscrit \textbf{corrigé} par un tiers (orthographe, typographie).
  \item \textbf{Format} choisi parmi ceux proposés par la ou les
        plateformes visées (un format commun facilite le multi-canal).
  \item \textbf{Mise en page} de l’intérieur (InDesign, Affinity
        Publisher, Scribus, \LaTeX{}, Vellum…) exportée en PDF imprimeur.
  \item \textbf{Couverture} réalisée sur le gabarit exact de la
        plateforme, avec prix, ISBN et code-barres.
  \item \textbf{Page de mentions légales} complète.
\end{checklist}

\subsection{Régler les formalités}

\begin{checklist}
  \item Demander vos \textbf{ISBN} à l’AFNIL (prévoir quelques jours à
        quelques semaines).
  \item Fixer le \textbf{prix public TTC} (loi Lang) en tenant compte des
        marges calculées au chapitre « Prix ».
  \item Préparer les informations \textbf{GPSR} (éditeur responsable
        dans l’UE, adresse postale et e-mail).
  \item Vérifier votre \textbf{statut} (artiste-auteur, micro-entreprise,
        société) et vos obligations déclaratives.
\end{checklist}

\subsection{Mettre en ligne et contrôler}

\begin{checklist}
  \item Créer la fiche du livre : \textbf{résumé} accrocheur,
        \textbf{catégories}, \textbf{mots-clés}, biographie.
  \item Téléverser les fichiers et corriger les alertes.
  \item Commander une \textbf{épreuve papier} et la relire en entier.
  \item Valider la publication, puis vérifier l’affichage du prix et de
        la disponibilité sur les principaux sites (comptez quelques jours
        à quelques semaines pour Fnac, Decitre, etc.).
  \item Effectuer le \textbf{dépôt légal} à la BnF.
\end{checklist}

\subsection{Faire connaître le livre}

L’IAD règle la fabrication et la logistique, pas la promotion. Les leviers
qui fonctionnent le mieux pour les petits éditeurs et les auteurs
indépendants :

\begin{itemize}
  \item une \textbf{newsletter} et une communauté (c’est le seul public
        qui vous appartient vraiment) ;
  \item les \textbf{librairies de proximité}, démarchées en personne,
        avec un dossier de présentation et des conditions claires
        (remise, dépôt-vente éventuel) ;
  \item les \textbf{salons} et rencontres, avec un petit stock imprimé
        à l’avance ;
  \item les \textbf{médiathèques}, souvent ouvertes aux auteurs locaux ;
  \item la \textbf{vidéo, les podcasts et les réseaux sociaux}, pour
        montrer les coulisses et l’humain derrière le livre.
\end{itemize}

\subsection{Les erreurs les plus fréquentes}

\begin{enumerate}
  \item Prendre l’ISBN gratuit d’une plateforme, puis vouloir changer de
        crèmerie.
  \item Fixer un prix trop bas, qui ne laisse aucune marge en librairie.
  \item Oublier le dépôt légal ou les mentions obligatoires.
  \item Valider sans commander d’épreuve papier.
  \item Pratiquer des prix différents selon les canaux (interdit par la
        loi Lang).
  \item Croire que « disponible en librairie » signifie « présent en
        rayon ».
\end{enumerate}
